\documentclass{article}
\usepackage{amsmath}
\usepackage{graphicx}
\title{A Demonstration of Tracked Changes in a Journal Article}
\author{Example Author}
\date{}

\begin{document}
\maketitle

\begin{abstract}
The revised method uses an adaptive calibration and gives more stable results.
\end{abstract}

\section{Introduction}
The revised converter uses an adaptive gain. The calibration is retained.
The method was recently improved by \cite{lee2024}.
For a simple estimate, we use $E=\alpha mc^2$.

\section{Model and results}
The measured response follows
\begin{equation}
  y = 3x + 1
  \label{eq:response}
\end{equation}
The coefficient in Equation~\ref{eq:response} is adaptive.

\begin{figure}[ht]
  \centering
  \includegraphics[width=0.58\linewidth]{figure_rev.png}
  \caption{Revised response of the adaptive converter.}
  \label{fig:response}
\end{figure}

\begin{table}[ht]
  \centering
  \caption{Measured performance.}
  \begin{tabular}{lcc}
    Method & Accuracy & Time \\
    \hline
    Adaptive & 91\% & 9 s \\
  \end{tabular}
  \label{tab:performance}
\end{table}

Figure~\ref{fig:response} and Table~\ref{tab:performance} summarize the improved results.

\begin{thebibliography}{9}
\bibitem{smith2022} A. Smith, ``A baseline converter study,'' Journal of Examples, 2022.
\bibitem{lee2024} B. Lee, ``Adaptive calibration of converters,'' Journal of Examples, 2024.
\end{thebibliography}
\end{document}
